\rtmtight
\begin{tblr}{width=\textwidth,colspec={|Q[c,m,2.1cm]|X[l,m]|},hlines,vlines,hspan=minimal,rowsep=1.5pt,colsep=3pt,row{1}={bg=headgray,font=\bfseries}}
\SetCell{c,m}\hfil\logo\hfil & \textbf{POLITEKNIK NEGERI MALANG}\par\textbf{JURUSAN TEKNOLOGI INFORMASI}\par\textbf{PROGRAM STUDI: D4 TEKNIK INFORMATIKA} \\
\SetCell[c=2]{l} \textbf{RENCANA TUGAS MAHASISWA (RTM) DAN RUBRIK PENILAIAN} & \\
Nama Mata Kuliah & Komputasi Hijau \\
Kode Mata Kuliah & RTI256108 \quad \textbf{sks / Jam perminggu:} 2 SKS/4 jam \quad \textbf{Semester:} 6 \\
Dosen Pengampu & Tim Pengajar Program Studi D4 TI \\
\end{tblr}
\assessmentblock{Tugas Analisis Dampak Lingkungan TI}{Case Method}{SCPMK0101-05201, SCPMK0509-05202}{Mahasiswa menganalisis dampak lingkungan dari infrastruktur TI nyata dan merancang strategi komputasi hijau menggunakan framework pengukuran emisi dan efisiensi energi.}{Mengidentifikasi infrastruktur TI target, mengukur atau mengestimasi carbon footprint, menganalisis konsumsi energi, merancang rekomendasi hemat energi, dan mendokumentasikan hasilnya.}{Laporan analisis carbon footprint, rencana strategi komputasi hijau, dan/atau bahan presentasi.}{Akurasi pengukuran dan analisis carbon footprint & 12 & Carbon footprint diestimasi dengan metodologi yang tepat (Scope 1/2/3), data konsumsi energi terdokumentasi lengkap, dan baseline ditetapkan dengan jelas. & Estimasi carbon footprint ada dan sebagian besar benar, tetapi metodologi atau data pendukung belum lengkap. & Pengukuran tidak menggunakan metodologi yang tepat atau data konsumsi energi tidak tersedia. \\ \\
Kelengkapan strategi rekomendasi komputasi hijau & 11 & Rekomendasi mencakup efisiensi hardware, virtualisasi, renewable energy, dan e-waste dengan target terukur dan rencana implementasi yang realistis. & Rekomendasi ada untuk sebagian aspek, tetapi target terukur atau rencana implementasi masih terbatas. & Rekomendasi bersifat umum, tidak spesifik, atau tidak dapat diimplementasikan pada konteks yang dianalisis. \\ \\
Kualitas laporan dan keterkaitannya dengan etika profesi TI & 7 & Laporan tersusun sistematis, mengaitkan temuan dengan tanggung jawab sosial dan etika profesi TI, serta menyajikan argumen yang kuat. & Laporan cukup lengkap tetapi kaitan dengan etika profesi atau tanggung jawab sosial masih lemah. & Laporan tidak terstruktur atau tidak mengaitkan analisis dengan dimensi etika dan tanggung jawab sosial. \\}

\assessmentblock{Kuis Konsep dan Regulasi Komputasi Hijau}{Kuis}{SCPMK0101-05201, SCPMK0509-05202}{Kuis tertulis untuk mengukur pemahaman konsep dasar komputasi hijau, regulasi keberlanjutan, dan prinsip efisiensi energi infrastruktur TI.}{Menjawab soal pilihan ganda dan/atau uraian terkait konsep green IT, carbon footprint, PUE, standar ISO 14001, dan Green Grid.}{Jawaban tertulis kuis.}{Pemahaman konsep komputasi hijau dan regulasi & 5 & Konsep green IT, carbon footprint, PUE, ISO 14001, dan Green Grid dijelaskan dengan benar dan contoh penerapannya tepat. & Sebagian besar konsep dipahami tetapi penjelasan atau penerapannya masih kurang lengkap atau ada kekeliruan minor. & Konsep tidak dipahami atau penjelasan menunjukkan miskonsepsi yang signifikan. \\ \\
Kemampuan menganalisis kasus efisiensi energi & 5 & Kasus dianalisis dengan tepat menggunakan metrik yang sesuai (PUE, kWh, emisi CO2), masalah diidentifikasi, dan solusi diusulkan secara logis. & Analisis kasus cukup tepat tetapi pemilihan metrik atau rekomendasi solusi belum sepenuhnya sesuai. & Analisis kasus tidak menggunakan metrik yang tepat atau tidak menunjukkan pemahaman efisiensi energi. \\}

\assessmentblock{UTS Carbon Footprint dan Efisiensi Energi}{UTS tertulis}{SCPMK0101-05201, SCPMK0509-05202}{Evaluasi tengah semester untuk mengukur kemampuan menganalisis carbon footprint dan merancang solusi efisiensi energi infrastruktur TI.}{Mengerjakan soal analisis kasus carbon footprint, perhitungan PUE, dan perancangan strategi efisiensi energi sesuai soal yang diberikan.}{Lembar jawaban UTS dengan perhitungan, analisis, dan rekomendasi tertulis.}{Ketepatan analisis carbon footprint dan perhitungan emisi & 8 & Perhitungan emisi CO2 menggunakan faktor emisi yang tepat, Scope 1/2/3 dibedakan dengan benar, dan interpretasi hasilnya akurat. & Perhitungan ada dan sebagian besar benar tetapi pembedaan scope atau interpretasi hasilnya masih kurang tepat. & Perhitungan salah atau tidak menggunakan faktor emisi yang sesuai standar. \\ \\
Kualitas perancangan strategi efisiensi energi & 7 & Strategi mencakup optimasi hardware, virtualisasi, dan pengelolaan beban kerja dengan estimasi penghematan energi yang realistis. & Strategi merancang beberapa aspek efisiensi tetapi estimasi penghematan atau cakupan solusinya masih terbatas. & Strategi tidak operasional atau tidak menunjukkan pemahaman tentang cara mengurangi konsumsi energi infrastruktur TI. \\}

\assessmentblock{PBL Audit Green IT}{Project Based Learning}{SCPMK0101-05201, SCPMK0509-05202}{Mahasiswa melakukan audit green IT komprehensif pada infrastruktur TI nyata atau simulasi, menghasilkan laporan audit dengan temuan dan rekomendasi keberlanjutan yang terukur.}{Merencanakan scope audit, mengumpulkan data konsumsi energi dan e-waste, menganalisis efisiensi infrastruktur, merancang roadmap keberlanjutan, dan mempresentasikan hasil audit.}{Laporan audit green IT, roadmap keberlanjutan, data pendukung, dan presentasi hasil.}{Kelengkapan dan kedalaman audit green IT & 9 & Audit mencakup energi (PUE, konsumsi server), e-waste, carbon footprint, dan kebijakan organisasi; temuan didukung data kuantitatif yang valid. & Audit mencakup sebagian besar aspek tetapi data kuantitatif untuk beberapa komponen masih kurang atau tidak diverifikasi. & Audit tidak sistematis, cakupan sangat terbatas, atau temuan tidak didukung data yang memadai. \\ \\
Kualitas rekomendasi dan roadmap keberlanjutan & 8 & Rekomendasi spesifik, terukur, dan diprioritaskan berdasarkan dampak dan feasibilitas; roadmap memiliki timeline dan KPI yang jelas. & Rekomendasi relevan tetapi prioritas, timeline, atau KPI roadmap masih perlu diperkuat. & Rekomendasi bersifat generik atau roadmap tidak operasional. \\ \\
Presentasi dan komunikasi temuan audit & 5 & Temuan dipresentasikan secara profesional, menggunakan visualisasi data yang tepat, dan mampu menjawab pertanyaan teknis secara meyakinkan. & Presentasi cukup jelas tetapi visualisasi atau kemampuan menjawab pertanyaan teknis masih terbatas. & Presentasi tidak terstruktur atau tidak mampu menyampaikan temuan audit secara jelas. \\}

\assessmentblock{UAS Desain Infrastruktur Berkelanjutan}{UAS presentasi dan defense}{SCPMK0509-05202}{Mahasiswa merancang infrastruktur TI berkelanjutan yang mengintegrasikan efisiensi energi, virtualisasi, renewable energy, dan prinsip green software engineering untuk konteks organisasi yang ditentukan.}{Menganalisis kebutuhan organisasi, merancang arsitektur infrastruktur hijau, menghitung estimasi pengurangan emisi dan penghematan energi, menyusun laporan teknis, dan mempresentasikan rancangan.}{Dokumen rancangan infrastruktur berkelanjutan, estimasi dampak lingkungan, laporan teknis, dan presentasi defense.}{Kelengkapan dan koherensi rancangan infrastruktur berkelanjutan & 10 & Rancangan mencakup semua komponen (server, pendingin, jaringan, storage, renewable energy) dengan justifikasi teknis yang kuat dan terintegrasi secara koheren. & Rancangan mencakup sebagian besar komponen tetapi integrasi atau justifikasi teknis untuk beberapa bagian masih kurang. & Rancangan tidak lengkap, komponen tidak terintegrasi, atau justifikasi teknis tidak memadai. \\ \\
Akurasi estimasi dampak lingkungan dan penghematan energi & 8 & Estimasi pengurangan emisi CO2 dan penghematan energi dihitung dengan metodologi yang tepat, target PUE realistis, dan proyeksi ROI keberlanjutan disajikan. & Estimasi ada dan sebagian besar menggunakan metodologi yang benar, tetapi beberapa kalkulasi atau asumsi perlu diperkuat. & Estimasi tidak menggunakan metodologi yang tepat atau angka yang disajikan tidak dapat dipertanggungjawabkan. \\ \\
Kualitas defense dan argumentasi teknis & 5 & Keputusan rancangan dipertahankan dengan argumen teknis, data referensi standar industri, dan analisis trade-off yang komprehensif. & Defense cukup meyakinkan tetapi beberapa keputusan tidak didukung argumen teknis atau referensi yang memadai. & Tidak mampu mempertahankan keputusan rancangan atau argumentasi tidak menunjukkan penguasaan materi. \\}

\begin{tblr}{width=\textwidth,colspec={|Q[l,m,3.2cm]|X[l,m]|},hlines,vlines,hspan=minimal,row{1-3}={bg=headgray},rowsep=1.5pt,colsep=3pt}
Jadwal Pelaksanaan: & Minggu 1--17 sesuai RPS. Setiap evaluasi memiliki RTM dan rubrik multi-kriteria. \\
Lain-Lain Yang Diperlukan: & LMS, tools pengukuran energi, simulator data center, laporan audit, dan perangkat presentasi. \\
Pustaka: & Mengacu pustaka utama dan pendukung pada bagian RPS. \\
\end{tblr}
